% Primitive: spectrum — authored from signal_processing/spectrum_manim.py,
% its README concepts table (seven scenes, watched in order) and plan 019's
% verified anchors (verifier report A–P plus the addendum Q–Z). Real
% sine-and-cosine probe pairs throughout — Steven W. Smith's "real DFT" —
% because the repo has built no complex number; the complex form is a
% forward pointer to Euler's formula and nothing here leans on it. The word
% is "detector": one frame gives one number per row. Every worked tone sits
% on a detector's frequency; two do not, both pointers — the computed
% 1500 Hz sine and the measured trumpet note. Problem answers
% are computed by answers/spectrum.py (the solve-gate anchor): exact in
% Q(√2), cross-checked against np.fft (whose imaginary part is the
% NEGATIVE of this chapter's plus-signed sine sum).
\chapter{The spectrum: a bank of detectors}

\begin{narrative}
A speech model never hears a sound. It reads a picture of one — a
spectrogram — and every column of that picture is the object this
chapter builds: a spectrum. The usual introduction makes the spectrum a
new kind of thing, reached through a new kind of number. This chapter
makes it an old one. Eight samples go in; five detectors listen to the
same eight samples, each at its own fixed frequency; and each detector
does nothing but the multiply-and-add that the \primref{random-variables}
chapter built for expectation. The spectrum is a column of weighted sums, read
off a bank of detectors — the formulation Steven W.~Smith calls the real
DFT~\cite{signal_processing-smith-guide-to-dsp-ch-8-the-dft}. All the
arithmetic runs on one grid: $N = 8$ samples at a sample rate of $sr =
8000$ per second.

\section{Pressure into numbers}

Sound is air pressure changing over time, and the simplest sound is a
pure tone. Picture it as a point going round a circle at a steady
speed; the microphone reads the point's \emph{height}. One lap per
millisecond is a thousand laps a second: a $1000$~Hz tone. A computer
cannot keep the whole motion. It keeps the height at regular instants —
$8000$ of them a second, so $8$ in this one millisecond, one every
$45^\circ$ of the lap. The $8$ stops on the circle \emph{are} the $8$
samples:
\begin{gather*}
  x[n] = \sin(2\pi n/8), \quad n = 0, \ldots, 7: \\
  \anchor{019.E.heights}.
\end{gather*}

\begin{center}
\begin{tikzpicture}[x=0.9cm, y=0.9cm]
  % PressureIntoNumbers: the 8 stops on the circle are the 8 samples —
  % stems, never a staircase (nothing between the samples is stored).
  \draw[Muted] (0, 0) circle (1.5);
  \foreach \n in {0,...,7}
    \fill[CoolInk] ({1.5*cos(45*\n)}, {1.5*sin(45*\n)}) circle (1.6pt);
  \draw[Muted] (2.6, 0) -- (10.6, 0);
  \draw[Muted, densely dashed] ({1.5*cos(45)}, {1.5*sin(45)}) -- (3.95, {1.5*sin(45)});
  \foreach [count=\i from 0] \v/\t in
      {0/0, 0.7071/0.71, 1/1, 0.7071/0.71, 0/0, -0.7071/-0.71, -1/-1, -0.7071/-0.71} {
    \draw[CoolInk, line width=1pt] (3+\i, 0) -- (3+\i, 1.5*\v);
    \fill[CoolInk] (3+\i, 1.5*\v) circle (2pt);
    \node[color=Muted] at (3+\i, -2.1) {\scriptsize $\i$};
    \node[color=CoolInk] at (3+\i, -2.6) {\small $\t$};
  }
  \node[left, color=Muted] at (2.45, -2.1) {\scriptsize sample $n$};
  \node[left, color=CoolInk] at (2.45, -2.6) {\small height};
  \node[color=Muted] at (0, 2.0) {\scriptsize $8$ stops, $45^\circ$ apart};
  \node[right, color=Muted] at (10.7, 0) {\scriptsize $1$ ms};
\end{tikzpicture}
\end{center}

The samples are drawn as stems because nothing between them is stored.
A staircase — each value held until the next — draws information the
recording does not contain; as Montgomery's demonstration puts it, the
stairsteps aren't really
there~\cite{signal_processing-montgomery-digital-show-and-tell-2013}.
Read the same point sideways and you get its \emph{shadow}: the same
motion a quarter-turn ahead,
\[
  \anchor{019.E.shadows}
\]
— two readings of one rotation. The height is the sine and the shadow
is the cosine, and the next section needs both.

Every digital recording starts here: $8000$ samples a second for
telephone speech, $16\,000$ for speech
models~\cite{signal_processing-jurafsky-martin-speech-and-language-processing-ch-15}.
One thing is planted for the last section: this tone gets $8$ samples
per lap; a faster tone gets fewer, and needs at least two. And real
sound is never one tone. Here are $8$ samples of a mix,
\[
  \anchor{019.R.samples},
\]
and the chapter's question: which tones are hiding in them?

\section{The probe}

To ask whether a $2000$~Hz tone is present, build one: a \emph{probe} is
a known tone that fits a whole number of laps into the window, and the
$2$-lap sine probe happens to be all integers at the $8$ stops — $0, 1,
0, -1, 0, 1, 0, -1$. Multiply the samples by the probe, stop by stop,
and add. That is the weighted sum $\sum_x x \cdot P(X = x)$ of the
\primref{random-variables} chapter, with two differences stated plainly: these
weights can be negative, and they add to $0$. They are a pattern to
match, not a probability.

\begin{center}
\begin{tikzpicture}[x=0.95cm, y=0.55cm]
  % TheProbe: match (every product a square, nothing can cancel) and
  % mismatch (four halves up, two wholes down) as product rows.
  \foreach [count=\i from 0] \a/\p in {0/0, 1/1, 0/0, -1/1, 0/0, 1/1, 0/0, -1/1} {
    \node[color=CoolInk] at (\i, 0) {\small $\a$};
    \node[color=Muted] at (\i, -1) {\small $\a$};
    \node[color=GoodInk] at (\i, -2.3) {\small $\p$};
  }
  \draw[Muted] (-0.45, -1.65) -- (7.45, -1.65);
  \node[left, color=CoolInk] at (-0.75, 0) {\small $2$-lap tone};
  \node[left, color=Muted] at (-0.75, -1) {\small $2$-lap probe};
  \node[left, color=GoodInk] at (-0.75, -2.3) {\small products};
  \node[right, color=AccentInk] at (7.75, -2.3) {\small sum $= 4$};
  \foreach [count=\i from 0] \a/\b/\p in
      {0/0/0, 0.71/0.71/0.5, 1/-1/-1, 0.71/0.71/0.5,
       0/0/0, -0.71/-0.71/0.5, -1/1/-1, -0.71/-0.71/0.5} {
    \node[color=CoolInk] at (\i, -4.6) {\small $\a$};
    \node[color=Muted] at (\i, -5.6) {\small $\b$};
    \node[color=WarmInk] at (\i, -6.9) {\small $\p$};
  }
  \draw[Muted] (-0.45, -6.25) -- (7.45, -6.25);
  \node[left, color=CoolInk] at (-0.75, -4.6) {\small $1$-lap tone};
  \node[left, color=Muted] at (-0.75, -5.6) {\small $3$-lap probe};
  \node[left, color=WarmInk] at (-0.75, -6.9) {\small products};
  \node[right, color=AccentInk] at (7.75, -6.9) {\small sum $= 0$};
\end{tikzpicture}
\end{center}

A \emph{match} first: the $2$-lap tone against the $2$-lap probe. Every
product is a square, so nothing can cancel: $\anchor{019.F.match}$. Then
a \emph{mismatch}: a $1$-lap tone against the $3$-lap probe. The
products take both signs — four halves up, two wholes down — and the
sum is exactly $0$. This detector hears nothing. It is Smith's
correlation picture, rebuilt small enough to
count~\cite{signal_processing-smith-ch-8-analysis-calculating-the-dft}.

One probe is not enough, though. Take a louder $2000$~Hz tone whose
samples run $\anchor{019.Q.tone}$ and repeat. The sine probe reads
$16$. Start the \emph{same tone} one sample later and it reads $12$:
the detector has been fooled by when the tone
starts~\cite{signal_processing-schaedler-seeing-circles-sines-and-signals}.
The repair is the shadow. Run the $2$-lap cosine probe as well and keep
the \emph{pair},
\begin{gather*}
  a_k = \sum_{n=0}^{7} x[n] \cos(2\pi k n/8), \qquad
  b_k = \sum_{n=0}^{7} x[n] \sin(2\pi k n/8), \\
  \text{reading} = \sqrt{a_k^2 + b_k^2} ,
\end{gather*}
both sums plus-signed. At the four starts the pair is
\[
  \anchor{019.Q.pairs} :
\]
one sample later is a quarter-turn of the pair, and its distance from
the origin is $\anchor{019.Q.distance}$ every time. The sine probe alone saw only the height of that turning point —
$\anchor{019.Q.sineonly}$.

\begin{center}
\begin{tikzpicture}[x=0.1cm, y=0.1cm]
  % TheProbe's pair plane: the pair-point walks a circle of radius 20.
  \draw[Muted, ->] (-27, 0) -- (27, 0);
  \draw[Muted, ->] (0, -27) -- (0, 27);
  \draw[Muted] (0, 0) circle (20);
  \draw[AccentInk, line width=1pt] (0, 0) -- (12, 16);
  \foreach \x/\y in {-16/12, -12/-16, 16/-12}
    \draw[Muted, densely dashed] (0, 0) -- (\x, \y);
  \foreach \x/\y in {12/16, -16/12, -12/-16, 16/-12}
    \fill[AccentInk] (\x, \y) circle (2pt);
  \node[above right, color=AccentInk] at (13, 17) {\small $(12, 16)$};
  \node[above left, color=AccentInk] at (-17, 13) {\small $(-16, 12)$};
  \node[below left, color=AccentInk] at (-13, -17) {\small $(-12, -16)$};
  \node[below right, color=AccentInk] at (17, -13) {\small $(16, -12)$};
  \node[right, color=Muted] at (28, 0) {\scriptsize cosine sum};
  \node[above, color=Muted] at (0, 28) {\scriptsize sine sum};
  \node[right, align=left, color=AccentInk] at (44, 8)
    {\small distance $20$\\ \small at every start};
  \node[right, align=left, color=Muted] at (44, -8)
    {\scriptsize one sample later:\\ \scriptsize a quarter-turn};
\end{tikzpicture}
\end{center}

\noindent The reading ignores when the tone starts — true for every
starting angle, shown here for the four the grid offers. And it measures
something the samples never show: $20 \div 4 = 5$ is the tone's
amplitude (the $4$ is the match sum above), while the largest sample is
only $4$. One detector, then, is a pair of probes: multiply, add, take
the distance. Every spectrum plot is this measurement repeated at each
frequency.

\section{The bank of detectors}

Build one detector for each whole number of laps the window can hold —
$0$, $1$, $2$, $3$, $4$ — and label the rows $0$, $1000$, $2000$,
$3000$ and $4000$~Hz \emph{before any sound arrives}. A row is an
address, not a discovery; what the literature calls a bin is a
detector's reading at a fixed address, never ``the frequency found in
the signal''~\cite{signal_processing-j-o-smith-mdft-spectral-bin-numbers}.
Feed the same $8$ samples to all five. A $2000$~Hz tone lights row $2$
alone, and it reads $4$. A $3000$~Hz tone started somewhere else
entirely — a cosine this time — lights row $3$ alone, again $4$. Then
the mix from the first section, which is Lyons' worked
example~\cite{signal_processing-lyons-understanding-dsp-3rd-ed-contents}
(the signal is his; the numbers here are computed): the five readings
are $\anchor{019.R.readings}$.

\begin{center}
\begin{tikzpicture}[x=1cm, y=0.8cm]
  % TheBankOfDetectors: the mix's 8 stems stay on screen beside the bank
  % (never a spectrum without its waveform); the end rows carry one probe.
  \draw[Muted] (-0.2, 0) -- (2.3, 0);
  \foreach [count=\i from 0] \v in
      {0.3536, 0.3536, 0.6464, 1.0607, 0.3536, -1.0607, -1.3536, -0.3536} {
    \draw[CoolInk, line width=0.8pt] (0.3*\i, 0) -- (0.3*\i, 0.9*\v);
    \fill[CoolInk] (0.3*\i, 0.9*\v) circle (1.4pt);
  }
  \node[color=CoolInk] at (1.05, -2.2) {\scriptsize the $8$ samples};
  \foreach [count=\k from 0] \hz/\probes/\len/\val in
      {0/{$c_0$}/0/0, 1000/{$c_1$ $s_1$}/2/4, 2000/{$c_2$ $s_2$}/1/2,
       3000/{$c_3$ $s_3$}/0/0, 4000/{$c_4$}/0/0} {
    \pgfmathsetmacro{\y}{2 - \k}
    \draw[Muted] (2.3, 0) -- (3.0, \y);
    \node[right] at (3.05, \y) {\scriptsize \hz~Hz};
    \node[draw=Muted, color=Muted, minimum width=1.15cm, minimum height=0.5cm,
      inner sep=2pt] (p\k) at (5.2, \y) {\scriptsize \probes};
    \node[draw=Muted, color=Muted, circle, inner sep=3.5pt] (s\k) at (6.5, \y)
      {\scriptsize $\Sigma$};
    \draw[Muted] (4.3, \y) -- (p\k.west);
    \draw[Muted] (p\k.east) -- (s\k.west);
    \draw[Muted] (s\k.east) -- (7.0, \y);
    \ifnum\len>0
      \fill[AccentInk] (7.1, \y-0.2) rectangle (7.1+\len, \y+0.2);
      \node[right, color=AccentInk] at (7.25+\len, \y) {\small $\val$};
    \else
      \fill[Muted] (7.1, \y-0.2) rectangle (7.14, \y+0.2);
      \node[right, color=Muted] at (7.3, \y) {\small $\val$};
    \fi
  }
  \node[right, color=AccentInk] at (9.75, 1) {\scriptsize pair $(0,\, 4)$};
  \node[right, color=AccentInk] at (8.75, 0) {\scriptsize pair $(1.41,\, -1.41)$};
\end{tikzpicture}
\end{center}

\noindent Two rows answer. Row $1$'s pair is $\anchor{019.R.rowonepair}$
and row $2$'s is $\anchor{019.R.rowtwopair}$ — the half-size tone
started at $135^\circ$, so neither of its probes alone tells the truth,
and its bar is the pair's distance, $2$: the previous section's lesson
inside the bank. Reading $\div\, 4$ gives the amplitudes, $1$ and
$0.5$. Turn the column of bars on its side and it is the plot everyone
calls \emph{the spectrum} — no new object, a column of weighted sums of
the same samples. The figure keeps the samples beside the bank on
purpose: students who are shown spectra without their waveforms
reliably fail to connect the two
pictures~\cite{signal_processing-wage-et-al-the-signals-and-systems-concept-inventory-2002}.

In code the bank already exists. \texttt{np.fft.rfft} returns these
five rows, each pair packed into one complex number — a device this
road has not built, so take it as a pointer: its real part is the
cosine sum and its imaginary part is \emph{minus} the sine sum,
$X[k] = a_k - i\,b_k$~\cite{signal_processing-j-o-smith-mathematics-of-the-dft-the-dft}.
Row $2$ of the mix, $\anchor{019.R.rowtwopair}$ here, is numpy's
$1.41 + 1.41i$. The reading — the distance — is the same either way.

\section{No double counting}

Why did the mix split so cleanly into $4$ and $2$, with nothing spilled
into the other rows? Two facts. First, each probe's sum on a mix is the
sum of its sums on the parts: a weighted sum of a sum is the sum of the
weighted sums, the same one-line reason the \primref{random-variables}
chapter gave for $E[X + Y] = E[X] + E[Y]$. (It is the \emph{pair} that
adds, cosine sum to cosine sum and sine sum to sine sum; the distance is
taken last, and distances do not add — the sums $12$ and $16$ made
$20$.) So it is enough to ask what each probe sums to on every other
probe. Second, the answer: nothing.

\begin{center}
\begin{tikzpicture}[x=0.72cm, y=0.55cm]
  % NoDoubleCounting: the 8 x 8 probe-against-probe table (anchor U).
  \foreach [count=\i from 0] \name/\d in
      {c_0/8, c_1/4, s_1/4, c_2/4, s_2/4, c_3/4, s_3/4, c_4/8} {
    \node[color=Muted] at (\i, 1.1) {\small $\name$};
    \node[color=Muted] at (-1.1, -\i) {\small $\name$};
    \foreach \j in {0,...,7} {
      \ifnum\i=\j
        \node[color=AccentInk] at (\j, -\i) {\small $\mathbf{\d}$};
      \else
        \node[color=Muted] at (\j, -\i) {\scriptsize $0$};
      \fi
    }
  }
  \node[right, align=left] at (8.3, -1.5)
    {\small $C(8, 2) = 28$ different pairs\\ \small all $28$ read exactly $0$};
  \node[right, align=left, color=Muted] at (8.3, -4.5)
    {\scriptsize $c$: cosine probe, $s$: sine probe\\
     \scriptsize the number: laps per window};
\end{tikzpicture}
\end{center}

There are $8$ probes — a cosine for each of the five rows, a sine for
each of the middle three — and so $64$ multiply-and-add sums. The
diagonal, each probe against itself, reads
$\anchor{019.U.diagonal}$. Off the diagonal there are
$\binom{8}{2} = \anchor{019.U.zeros}$ different pairs, counted the way
the \primref{counting-rules} chapter counts them, and every one reads exactly $0$.
That is \emph{computed, not proved}: an exhaustive check at $N = 8$. The
reason it holds for every $N$ needs a tool this road has not built —
Euler's formula — and waits for
it~\cite{signal_processing-j-o-smith-mdft-orthogonality-of-the-dft-sinusoids}.
The consequence needs no waiting: each detector is exactly deaf to the
other detectors' tones, so no bar can borrow from another.

Why $4$ on the diagonal? Squares never cancel, and at every stop
height$^2$ + shadow$^2 = 1$ — it is the circle's radius, by Pythagoras.
The squared sine probe and the squared cosine probe run
\begin{gather*}
  \anchor{019.V.sinesquares} \\
  \anchor{019.V.cosinesquares}
\end{gather*}
— the same values in a different order, summing to $1$ at each stop. Eight stops share out $8$
between two probes that take the same values — $4$
each~\cite{signal_processing-j-o-smith-mdft-norm-of-the-dft-sinusoids}.

Why $8$ at the two ends? At $0$ laps and at $4$ laps the sine probe is
all zeros — $\sin 0$ and $\sin(\pi n)$ vanish at every
stop~\cite{signal_processing-smith-ch-8-dft-basis-functions} — so the
cosine has no partner and keeps the whole $8$. The end rows are
therefore different in two ways, stated once: they read $\div\, 8$, not
$\div\, 4$~\cite{signal_processing-smith-ch-8-synthesis}; and with only
one probe they \emph{do} care when the tone starts. A full-size tone on
an end row reads $\anchor{019.B.ends}$ as its start moves through
$0^\circ$, $30^\circ$, $45^\circ$ and $90^\circ$. This is why every
worked tone in this chapter lives on rows $1$ to $3$.

The bookkeeping closes the argument: $1 + 2 \cdot 3 + 1 = 8$ readings
from $8$ samples~\cite{signal_processing-smith-ch-8-notation-and-format-of-the-real-dft}.
Nothing is lost, which is why the transform can be undone — the $8$
readings rebuild the $8$ samples. That is a promise, not something shown
here. And one warning rides along: each detector is deaf at the
\emph{other detectors'} frequencies — and only
there~\cite{signal_processing-j-o-smith-mdft-frequencies-in-the-cracks,signal_processing-lyons-using-the-dft-as-a-filter-2013}.
The next section meets a tone that sits somewhere else.

\section{What sets the spacing}

What decides where the detectors listen? Probe $k$ fits $k$ laps into
the window, and the window lasts $N \div sr$ seconds — here $8 \div
8000$~s, one millisecond. So $k$ laps per millisecond is $k \times
1000$~Hz, and in general
\[
  f_k = k \cdot \frac{sr}{N}, \qquad
  \text{spacing} = \frac{sr}{N} = \frac{1}{\text{duration}} .
\]
The spacing between detectors is one over how long the window listens.
Three banks, side by side, separate the two things a designer can buy:

\begin{center}
\begin{tabular}{lll}
  $sr = 8000$, $N = 8$ & $\anchor{019.X.first}$ & rows to $4000$~Hz \\
  $sr = 16\,000$, $N = 16$ & $\anchor{019.X.second}$ & rows to $8000$~Hz \\
  $sr = 8000$, $N = 16$ & $\anchor{019.X.third}$ & rows to $4000$~Hz \\
\end{tabular}
\end{center}

\noindent Sampling twice as fast for the same millisecond leaves the
rows $1000$~Hz apart — the faster rate bought \emph{reach}, not detail.
Listening twice as long put them $500$~Hz apart: only duration buys
spacing. The row count is $N/2 + 1$ in every case.

This is a speech model's front end. Whisper takes $400$ samples at
$16\,000$~Hz — a window of $\anchor{019.L.window}$~ms — which is
$\anchor{019.L.rows}$ rows, $\anchor{019.L.spacing}$~Hz apart, up to
$8000$~Hz~\cite{signal_processing-radford-et-al-robust-speech-recognition-2022,signal_processing-openai-whisper-audio-py-at-8609812}.
And the spacing has a cost, which this chapter shows and does not
explain. A $1500$~Hz sine starting at $0$ fits one and a half laps into
the window — its $8$ samples show it: not a whole number of laps — so it
sits between two rows of the small bank, and \emph{every} row answers —
on rows $0$ to $4$,
\[
  \anchor{019.T.teaser} .
\]
That smear is
\emph{leakage}, and it is the windowing chapter's subject — the one
computed tone in this chapter that does not sit on a detector's
frequency (the next section's trumpet note is the other, measured
rather than computed). Two more
pointers along the same road: slide the window along the recording and
each detector's output becomes a signal in its own right, which is when
the bank earns the name
\emph{filterbank}~\cite{signal_processing-j-o-smith-spectral-audio-signal-processing-dft-filter-bank};
and a \emph{mel} filterbank will regroup Whisper's
$\anchor{019.L.rows}$ rows into $80$.

\section{The spectrum in decibels}

Add a third tone to the mix, a hundred times quieter than the first:
amplitude $0.01$ at $3000$~Hz. The readings are
$\anchor{019.S.readings}$, and on ordinary bars the third tone is
invisible. The cure is the counting strip of the \primref{logarithms} chapter.
There is one definition: a level in decibels is ten times the common
logarithm of a \emph{power}
ratio~\cite{signal_processing-martin-decibel-the-name-for-the-transmission-unit-1929}.
Power goes as amplitude squared, and the law $\log(xy) = \log x + \log
y$ — multiplying is adding — turns the square into a factor of $2$ out
front~\cite{signal_processing-j-o-smith-mdft-decibels,signal_processing-nist-sp-811-8-7-logarithmic-quantities-and-units}:
\[
  L = 10 \log_{10} \frac{P}{P_0} = 20 \log_{10} \frac{A}{A_0}
  \ \text{dB} .
\]
A decibel figure means nothing without its reference; here it is the
loudest row. That row is ratio $1$, and $\log 1 = 0$~dB. Half the
reading is $20 \log_{10} \tfrac12 = -20 \times \anchor{005.log10two} =
\anchor{019.S.halfdb}$~dB, on the logarithm chapter's own value of
$\log_{10} 2$. A hundredth is two tenfoldings down:
$\anchor{019.S.hundredthdb}$~dB — and now the third tone has a bar.

\begin{center}
\begin{tikzpicture}[x=1cm, y=1cm]
  % TheSpectrumInDecibels: the same three readings on linear bars and on a
  % dB axis whose floor is labelled (a bar's height is a claim).
  \draw[Muted] (-0.5, 0) -- (3.9, 0);
  \foreach [count=\k from 0] \h/\t/\hz in
      {0/0/0, 2/4/1000, 1/2/2000, 0.02/0.04/3000, 0/0/4000} {
    \ifdim\h pt>0pt
      \fill[AccentInk] (0.85*\k-0.18, 0) rectangle (0.85*\k+0.18, \h);
      \node[above, color=AccentInk] at (0.85*\k, \h+0.05) {\scriptsize $\t$};
    \else
      \node[above, color=Muted] at (0.85*\k, 0.05) {\scriptsize $\t$};
    \fi
    \node[below, color=Muted] at (0.85*\k, -0.12) {\tiny \hz};
  }
  \node[color=Muted] at (1.7, 3.2) {\scriptsize readings, by row (Hz)};
  \draw[Muted] (5.6, 0) -- (11.0, 0);
  \foreach [count=\k from 0] \h/\t/\hz in
      {0/{-\infty}/0, 2.5/0/1000, 2.2/{-6.02}/2000, 0.5/{-40}/3000, 0/{-\infty}/4000} {
    \ifdim\h pt>0pt
      \fill[AccentInk] (6.2+1.1*\k-0.18, 0) rectangle (6.2+1.1*\k+0.18, \h);
      \node[above, color=AccentInk] at (6.2+1.1*\k, \h+0.05) {\scriptsize $\t$};
    \else
      \node[above, color=Muted] at (6.2+1.1*\k, 0.05) {\scriptsize $\t$};
    \fi
    \node[below, color=Muted] at (6.2+1.1*\k, -0.12) {\tiny \hz};
  }
  \node[left, color=Muted] at (5.5, 0) {\scriptsize $-50$ dB};
  \node[color=Muted] at (8.4, 3.2) {\scriptsize dB re: the loudest row, by row (Hz)};
\end{tikzpicture}
\end{center}

The empty rows are the boundary reading the logarithm chapter flagged:
$\log 0 = -\infty$. So $0$~dB means \emph{equal to the reference} — it
is never silence; silence is $-\infty$~dB. Real systems therefore pick
a floor. Whisper keeps $8$ decades of power below its loudest value —
$\anchor{019.Z.floor}$~dB — and clamps the
rest~\cite{signal_processing-openai-whisper-audio-py-at-8609812}. Two
conversions are worth carrying: doubling an amplitude is
$\anchor{019.I.amplitudedouble}$~dB, doubling a power
$\anchor{019.I.powerdouble}$~dB.

A real note shows what the axis is for. Take $4096$ samples of a
trumpet at $22\,050$ per second: $\anchor{019.M.rows}$ detectors,
$\anchor{019.M.spacing}$~Hz apart. Its first five harmonics, in Hz, and their
levels in dB re the loudest:
\begin{gather*}
  \anchor{019.M.harmonics} \\
  \anchor{019.M.levels}
\end{gather*}
The \emph{second} harmonic is the loudest, the fundamental $1.56$~dB
below it. (The recording is Mihai Sorohan's, from
Freesound, CC~BY~4.0, measured from librosa's converted
copy~\cite{signal_processing-sorohan-solo-trumpet-06in-f-90bpm-wav-freesound,signal_processing-librosa-example-files};
the plot it rebuilds is the Hugging Face audio course's~\cite{signal_processing-hugging-face-audio-course-introduction-to-audio-data}.)
One honest pointer: a real note lands between rows, so a window shaped
those levels — the windowing chapter again.

The same ruler is everywhere. Decibels are $10 \times \log_{10}$ of a
power ratio; pH is $-\log_{10}$ of hydrogen-ion
activity~\cite{signal_processing-buck-et-al-measurement-of-ph-iupac-2002};
stellar magnitude is $-2.5 \times \log_{10}$ of a brightness ratio, so
that five magnitudes are a factor of
$100$~\cite{signal_processing-pogson-magnitudes-of-thirty-six-of-the-minor-planets-1856};
semitones are $12 \times \log_2$ of a frequency ratio, which makes a
$3{:}2$ fifth $\anchor{019.Y.fifth}$ of them. Each is a constant times
the log of a ratio — the logarithm chapter's strip with its own stride.
That the ear itself hears in logs is an approximate motivation, not a
law~\cite{signal_processing-wikipedia-weber-fechner-law}.

\section{The fold at half the sample rate}

Why does the bank stop at row $4$ — at $4000$~Hz, half the sample rate?
Because a row $5$ would be row $3$ again. At the $8$ stops the $7$-lap
cosine probe \emph{is} the $1$-lap cosine probe, number for number, and
the $7$-lap sine probe is the $1$-lap one with every sign flipped:
\[
  \cos\bigl(2\pi (8 - k) n/8\bigr) = \cos(2\pi k n/8), \qquad
  \sin\bigl(2\pi (8 - k) n/8\bigr) = -\sin(2\pi k n/8) .
\]
Likewise $6$ laps is $2$, and $5$ laps is $3$. Run the table of zeros
past row $4$ and it lights up: $c_3$ against $c_5$ adds to
$\anchor{019.U.cosfold}$, $s_3$ against $s_5$ to
$\anchor{019.U.sinfold}$ — not $0$. A detector above row $4$ would only
repeat one below it.

The reason is on the first section's circle. A $7000$~Hz tone steps
$7/8$ of a lap between samples, visiting the angles (in degrees)
\[
  \anchor{019.W.angles}
\]
— exactly the stops of a point stepping $1/8$ of a lap \emph{backward}. Same shadow, height flipped: it is the
wagon wheel that seems to turn the wrong
way~\cite{signal_processing-wikipedia-aliasing}. Two cosines, $1000$~Hz
and $7000$~Hz, pass through the same $8$ samples, the samples cannot
say which was played, and it is the $1000$~Hz row that fires. On the
frequency ruler, sampled at $8000$~Hz:
\[
  \anchor{019.W.ruler} .
\]

\begin{center}
\begin{tikzpicture}[x=1.6cm, y=1cm]
  % TheFoldAtNyquist: the frequency ruler folded at sr/2 — 5 -> 3,
  % 6 -> 2, 7 -> 1 only (past 8000 Hz the pattern changes, anchor W).
  \draw[Muted, line width=1.2pt] (0, 0) -- (4, 0);
  \foreach [count=\k from 0] \hz in {0, 1000, 2000, 3000, 4000} {
    \draw[Muted] (\k, 0.1) -- (\k, -0.1);
    \node[below] at (\k, -0.2) {\small \hz};
  }
  \draw[WarmInk, densely dashed, line width=1pt]
    (4, 0) arc (-90:90:0.18 and 0.6) -- (0.6, 1.2);
  \foreach \f/\x in {5000/3, 6000/2, 7000/1} {
    \node[above, color=WarmInk] at (\x, 1.35) {\small \f};
    \draw[WarmInk, ->, line width=0.8pt] (\x, 1.12) -- (\x, 0.22);
  }
  \node[right, align=left, color=AccentInk] at (4.4, 0.6)
    {\small the fold:\\ \small $sr/2 = 4000$ Hz};
  \node[below, color=Muted] at (4.75, -0.2) {\small Hz};
\end{tikzpicture}
\end{center}

\noindent (The ruler stops at $7000$ on purpose. From $8000$~Hz up the
picture is no longer a mirror, and ``cosine the same, sine flipped''
stops being the rule.)

Now the first section's promise is cashed. It said a tone needs at least
two samples per lap; strictly, it needs \emph{more than} two: it must lie
\emph{below} half the
sample rate ($sr/2$), not at
it~\cite{signal_processing-j-o-smith-mdft-sampling-theorem}. Exactly at
$sr/2$ a sine lands on $0$ at every stop, and the samples are silence —
the end row's missing partner, met from the other side. (Textbooks
disagree about what ``the Nyquist frequency'' names, so this chapter
says ``half the sample rate''
instead~\cite{signal_processing-oppenheim-schafer-discrete-time-signal-processing-3rd-ed}.)
And the impostor is a \emph{clean lower tone}, not noise: nothing
downstream can tell it from the real thing, so whatever lies above
$sr/2$ has to be removed \emph{before} sampling. Hence $8000$~Hz for
telephone speech and $16\,000$~Hz for speech models — where, by the same
fold, $\anchor{019.H.speechrate}$~Hz. In numpy the fold is visible in
the output itself: the upper half of \texttt{fft}'s output is this
mirror, and \texttt{rfft} returns the bank. The sampling theorem's credit
line: E.~T.~Whittaker 1915, Kotelnikov 1933, Shannon
1949~\cite{signal_processing-shannon-communication-in-the-presence-of-noise-1949,signal_processing-l-ke-the-origins-of-the-sampling-theorem-1999}.

\paragraph{Where this goes.} The series closes on a map. \emph{Which
tones are in it?} — the bank of detectors. \emph{Quiet beside loud?} —
decibels. Both are built. \emph{Echoes and smoothing?} —
convolution. \emph{A tone between rows?} — windowing. \emph{Tones that
change over time?} — the short-time transform, whose columns stacked
side by side are the spectrogram. \emph{The ear's own grouping?} — mel.
Those are the road ahead, convolution first because leakage is best
explained by it. Three further promises leave
this chapter. The complex form — one complex number per row instead of
a pair, and the reason the $\anchor{019.U.zeros}$ zeros hold for every
$N$ — waits on Euler's formula. The inverse, rebuilding the samples
from the readings, is unbuilt. And the fast Fourier transform computes
these same readings in far fewer steps by divide and conquer — an
algorithms story, Cooley and Tukey's in
1965~\cite{signal_processing-cooley-tukey-an-algorithm-for-the-machine-calculation-1965},
though Gauss had the general formula around 1805, published in
1866~\cite{signal_processing-heideman-johnson-burrus-gauss-and-the-history-of-the-fft}.
\end{narrative}

\section*{Practice problems}

\begin{problem}
Eight samples are taken at $8000$ per second. Here they are, above the
$2000$~Hz detector's two probes at the same $8$ stops:
\[
\begin{array}{lrrrrrrrr}
  \text{samples} & 5 & 12 & -5 & -12 & 5 & 12 & -5 & -12 \\
  \text{cosine probe} & 1 & 0 & -1 & 0 & 1 & 0 & -1 & 0 \\
  \text{sine probe} & 0 & 1 & 0 & -1 & 0 & 1 & 0 & -1
\end{array}
\]
Compute the detector's pair (cosine sum, sine sum), its reading, and
the amplitude of the tone; compare the amplitude with the largest
sample. Then give the pair of the $1000$~Hz detector, whose probes are
$\cos(2\pi n/8)$ and $\sin(2\pi n/8)$ for $n = 0, \ldots, 7$.
\begin{hint}
Each sum has only four nonzero terms. For the $1000$~Hz detector, pair
each sample with the one four places later.
\end{hint}
\begin{solution}
Cosine sum: $5 + 5 + 5 + 5 = 20$ (the $-5$s meet the $-1$s). Sine sum:
$12 + 12 + 12 + 12 = 48$. The pair is $(20, 48)$ and the reading is
$\sqrt{20^2 + 48^2} = 52$. On an interior row the amplitude is the
reading $\div\, 4$: $13$ — larger than any sample, the largest being
$12$, because the tone is $5 \cos + 12 \sin$ and peaks between stops
($5$–$12$–$13$). For the $1000$~Hz detector, the samples satisfy $x[n +
4] = x[n]$ while both $1$-lap probes change sign four places later, so
the terms cancel in pairs: $(0, 0)$.
\end{solution}
\end{problem}

\begin{problem}
All signals here are $8$ samples at $8000$ per second, read by the
five-detector bank ($0$, $1000$, $2000$, $3000$, $4000$~Hz). (a) The
signal is $2.5 \sin(2\pi \cdot 3000\, t + 50^\circ)$. Give all five
readings. (b) The signal is $3 \sin(2\pi \cdot 1000\, t) + 0.5
\cos(2\pi \cdot 3000\, t)$. Give all five readings, and name the two
facts that make the other rows read exactly $0$.
\begin{hint}
On rows $1$ to $3$ a tone of amplitude $A$ sitting on the row's
frequency reads $4A$, whatever its starting angle.
\end{hint}
\begin{solution}
(a) Only row $3$ answers: $4 \times 2.5 = 10$, so the readings are $0,
0, 0, 10, 0$; the $50^\circ$ turns the pair, not its distance. (b) $0,
12, 0, 2, 0$: row $1$ reads $4 \times 3$, row $3$ reads $4 \times 0.5$.
The two facts: each probe's sum on a mix is the sum of its sums on the
parts (a weighted sum is linear; the pair adds, and the distance is
taken last), and every probe sums to exactly $0$ against every other
probe — all $28$ pairs at $N = 8$ — so each tone reaches only its own
row.
\end{solution}
\end{problem}

\begin{problem}
For each bank give the window's duration, the spacing between
detectors, the number of rows, and the highest row's frequency. (a)
Whisper's: $sr = 16\,000$~Hz, $N = 400$. (b) $sr = 48\,000$~Hz, $N =
960$. (c) You need detectors $10$~Hz apart at $sr = 16\,000$~Hz: how
many samples, and how long a window? If you switch to $sr =
48\,000$~Hz and keep that $N$, what is the spacing — and what $N$
restores $10$~Hz?
\begin{hint}
Spacing $= sr/N = 1/\text{duration}$; rows $= N/2 + 1$; the top row is
$sr/2$.
\end{hint}
\begin{solution}
(a) $400/16\,000$~s $= 25$~ms; $16\,000/400 = 40$~Hz; $201$ rows; up to
$8000$~Hz. (b) $20$~ms; $50$~Hz; $481$ rows; up to $24\,000$~Hz. (c)
$N = 16\,000/10 = 1600$ samples, $100$~ms ($801$ rows). At $48\,000$~Hz
the same $N$ lasts a third as long and the spacing is $48\,000/1600 =
30$~Hz — coarser. $10$~Hz needs $N = 4800$, which is again $100$~ms:
spacing is bought with time, at any sample rate.
\end{solution}
\end{problem}

\begin{problem}
Use $\log_{10} 2 \approx 0.301$. Levels are in dB relative to the
loudest row unless a power ratio is named. (a) The loudest row reads
$4$ and another reads $0.5$: what is the second row's level? (b) A row
sits at $-12.04$~dB: what is its amplitude ratio to the loudest row,
and the power ratio? (c) What is the level of a tenth of the amplitude?
Of a tenth of the power? (d) A tone has half the \emph{power} of the
reference: what is its level, and its amplitude ratio? (e) The loudest
row reads $4$; another row sits at $-26.02$~dB. What does it read?
\begin{hint}
$20 \log_{10}$ of an amplitude ratio, $10 \log_{10}$ of a power ratio;
split every ratio into halvings and tenfoldings.
\end{hint}
\begin{solution}
(a) $0.5/4 = 1/8$, three halvings: $-20 \times 3 \times 0.301 =
-18.06$~dB. (b) $-12.04 = -20 \times 2 \times 0.301$, two halvings: an
amplitude ratio of $1/4$, a power ratio of $1/16$. (c) $20 \log_{10}
\tfrac{1}{10} = -20$~dB; for power, $10 \log_{10} \tfrac{1}{10} =
-10$~dB. (d) $10 \log_{10} \tfrac12 = -3.01$~dB; amplitude goes as the
square root of power, so the ratio is $1/\sqrt2 \approx 0.707$. (e)
$-26.02 = -20 - 6.02$: a tenth, then a half, so the ratio is $1/20$ and
the row reads $4/20 = 0.2$.
\end{solution}
\end{problem}

\begin{problem}
(a) A $5000$~Hz tone is sampled at $8000$ per second. At what frequency
is it heard? (b) At $16\,000$ per second, where are tones of $5000$,
$9000$ and $10\,000$~Hz heard? (c) Write the $8$ samples of $\sin(2\pi
\cdot 7000\, t)$ taken at $8000$ per second, starting at $t = 0$, to
two decimals. Which $1000$~Hz tone has exactly these samples? Give the
$1000$~Hz detector's pair (cosine sum, sine sum) and reading. (d) The
same for $\cos(2\pi \cdot 7000\, t)$: pair and reading on the
$1000$~Hz detector.
\begin{hint}
Between $sr/2$ and $sr$ a tone lands at $sr - f$, its cosine part
unchanged and its sine part negated.
\end{hint}
\begin{solution}
(a) $8000 - 5000 = 3000$~Hz. (b) $5000$~Hz is below $sr/2 = 8000$~Hz
and stays put; $9000 \to 7000$~Hz and $10\,000 \to 6000$~Hz. (c) The samples are
\[
  0,\ -0.71,\ -1,\ -0.71,\ 0,\ 0.71,\ 1,\ 0.71 ,
\]
those of $-\sin(2\pi \cdot 1000\, t)$: a $1000$~Hz sine turned upside
down. The pair is $(0, -4)$
and the reading $4$. (d) The samples equal those of $\cos(2\pi \cdot
1000\, t)$ exactly: pair $(4, 0)$, reading $4$. Either way the
impostor is a clean $1000$~Hz tone at full reading.
\end{solution}
\end{problem}

\begin{problem}
On the $N = 8$, $sr = 8000$~Hz grid the top row listens at $4000$~Hz,
exactly $sr/2$. (a) Write the samples of $\cos(2\pi \cdot 4000\, t)$.
What does row $4$ read, and what divisor recovers the amplitude $1$?
(b) Write the samples of $\sin(2\pi \cdot 4000\, t)$. What does every
row read? (c) Write the samples of $3 \cos(2\pi \cdot 4000\, t +
60^\circ)$. What does row $4$ read, and what amplitude does the
end-row divisor report?
\begin{hint}
At $4$ laps the sine probe is all zeros, so row $4$ is its cosine sum
alone; $\cos(\pi n + \varphi) = (-1)^n \cos\varphi$.
\end{hint}
\begin{solution}
(a) $1, -1, 1, -1, 1, -1, 1, -1$. Row $4$'s cosine probe is the same
list, so it reads $8$; the end rows divide by $8$, not $4$: $8/8 = 1$.
(b) $\sin(\pi n) = 0$ at every stop: the samples are silence and every
row reads $0$ — which is why a tone must lie strictly below $sr/2$.
(c) $3 \cos 60^\circ = 1.5$, so the samples are $1.5, -1.5, 1.5,
\ldots$ and row $4$ reads $8 \times 1.5 = 12$; dividing by $8$ reports
$1.5$, half the true amplitude. With no sine partner the end row reads
$A \cdot N \cdot |\cos\varphi|$ — it depends on when the tone starts.
\end{solution}
\end{problem}

\begin{problem}
Exactly one of these five claims is true. Which — and what is wrong
with each of the others? (a) ``Between two samples a recording holds
the earlier value, so the honest picture of $8$ samples is a
staircase.'' (b) ``$8$ samples at $8000$ per second put the detectors
$1000$~Hz apart; $16$ samples at $16\,000$ per second put them $500$~Hz
apart.'' (c) ``A row at $0$~dB is silent.'' (d) ``A $7000$~Hz tone
sampled at $8000$ per second shows up as noise spread across the
rows.'' (e) ``Starting a $2000$~Hz tone a quarter of a lap later
changes the $2000$~Hz detector's pair but not its reading.''
\begin{hint}
Check (b) with $sr/N$, (c) with $\log 1$ and $\log 0$, (d) with the
fold.
\end{hint}
\begin{solution}
Only (e) is true: the pair turns by a quarter-turn and its distance
stays put. (a) False: nothing between the samples is stored — a sample
is a value at an instant, drawn as a stem. (b) False: both windows last
$1$~ms, so both banks are $1000$~Hz apart; the faster rate bought reach
(rows to $8000$~Hz), and $500$~Hz takes a $2$~ms window. (c) False:
$0$~dB means equal to the reference ($\log 1 = 0$); silence is
$-\infty$~dB. (d) False: it lands on $1000$~Hz as a clean tone — row
$1$ reads $4$ and every other row $0$.
\end{solution}
\end{problem}

\begin{problem}[stretch]
The tone of the first problem starts one sample later, so the samples
read $-12, 5, 12, -5, -12, 5, 12, -5$; two samples later they read $-5,
-12, 5, 12, \ldots$; three samples later $12, -5, -12, 5, \ldots$.
Give the $2000$~Hz detector's pair (cosine sum, sine sum) and its
reading at each of the three delays, and say what the sine probe alone
would have reported at delays $0$ to $3$. Through what angle does one
sample's delay turn the pair — and through what angle would it turn a
$1000$~Hz tone's pair on the $1000$~Hz detector?
\begin{hint}
The probes have not moved; only the samples have. A tone of $f$~Hz
advances $f/sr$ of a lap per sample.
\end{hint}
\begin{solution}
Delay $1$: $(-48, 20)$. Delay $2$: $(-20, -48)$. Delay $3$: $(48,
-20)$. Each is the previous pair turned a quarter-turn anticlockwise —
$(a, b) \mapsto (-b, a)$ — and the reading is $52$ every time. The sine
probe alone reports $48, 20, -48, -20$: the same tone, four different
answers. One sample is $2000/8000$ of a lap, $90^\circ$; for a
$1000$~Hz tone it is $1000/8000$ of a lap, $45^\circ$.
\end{solution}
\end{problem}
